\section{How to Run}

\subsection{Native (WSL / Linux) --- Not recommended}
The configurations can be a little messy because requires the SICStus interpreter and the DALI logic, but here are presente a step-by-step guide:
\begin{enumerate}
  \item Clone the DALI repo from \url{https://github.com/AAAI-DISIM-UnivAQ/DALI} and navigate to the \texttt{DALI/Examples} folder.
  \item Clone the GitHub repo project from \url{https://github.com/Mik1810/f1_dali_project} under \texttt{Examples} folder.
\end{enumerate}
Then, execute the following commands in the terminal:\\

\begin{lstlisting}[language=bash, numbers=none]
cd DALI/Examples/f1_race
bash startmas.sh
\end{lstlisting}

\texttt{startmas.sh} performs the following steps:
\begin{enumerate}[leftmargin=2em]
  \item Kills any stale SICStus / tmux processes on port 3010.
  \item Runs \texttt{generate\_agents.py} to sync DALI files with \texttt{agents.json}.
  \item Launches all agents inside a single \texttt{tmux}
        session named \texttt{f1\_race}.
\end{enumerate}

\subsection{Web Dashboard --- Recommended}

If you want to use the web dashboard, which provides real-time 
logs and an animated circuit view, run the following commands
 in terminal:\\

\begin{lstlisting}[language=bash, numbers=none]
cd f1_race
bash ui/run.sh
\end{lstlisting}

\noindent
\texttt{run.sh} creates a local Python venv (\texttt{ui/.venv}) on the first run and
installs Flask automatically. Subsequent launches or skip \texttt{pip install} via a stamp
file (\texttt{.venv/.installed\_stamp}) and start it.

Open \url{http://localhost:5000} in a browser. The dashboard shows all agent panes
side-by-side with real-time auto-scrolling. If \texttt{agents.json} is modified while the
dashboard is open, it auto-detects the change within 5\,seconds and rebuilds the grid
without a page reload.

\begin{figure}[H]
  \centering
  \includegraphics[width=\textwidth]{../../imgs/07_consoles.png}
  \caption{Web dashboard --- Logs tab: agent panes with real-time output, control buttons and command bar.}
  \label{fig:ui_consoles}
\end{figure}

\subsubsection{Dashboard Controls}

\begin{longtable}{@{} l p{9cm} @{}}
\toprule
\textbf{UI element} & \textbf{Function} \\
\midrule
\endhead
\textbf{Restart MAS}    & Kills SICStus + tmux session, reruns \texttt{startmas.sh} \\
\textbf{Deploy SC}      & Sends deploy message to safety car immediately \\
\textbf{Recall SC}      & Recalls the safety car \\
Agent / Command bar     & Sends an arbitrary Prolog command to any agent pane \\
Pin button ($\downarrow$) & Toggles auto-scroll for that pane \\
Clear button ($\times$)  & Clears pane output \\
Minimise button ($-$)   & Minimises pane to tray \\
\bottomrule
\end{longtable}

\begin{figure}[H]
  \centering
  \includegraphics[width=\textwidth]{../../imgs/08_podium.png}
  \caption{Final podium modal displayed at race end, showing driver positions and total accumulated times.}
  \label{fig:ui_podium}
\end{figure}

\subsubsection{Circuit Tab}

The dashboard exposes two tabs selectable via the tab bar at the top of the page:

\begin{itemize}[leftmargin=2em]
  \item \textbf{Logs}: the default view showing all agent panes side-by-side with
        real-time output.
  \item \textbf{Circuit}:  an animated canvas visualisation of the race circuit.
\end{itemize}

\noindent
Clicking a tab label switches the view instantly; no page reload is required. The
\texttt{Logs} tab remains fully active in the background (polling continues) while
the \texttt{Circuit} tab is open, so no agent output is ever lost.

The \textbf{Circuit tab} renders a track on an canva and
animates each car along it in proportional time. Key features include:

\begin{itemize}[leftmargin=2em]
  \item \textbf{Start lights sequence}:  five red lights illuminate one per second
        before going out at race start, matching the semaphore agent's timing.
  \item \textbf{Car tokens}: each car is drawn as a coloured dot with its team
        border glow, a lap-count badge, and a driver label. Position on track reflects
        the real fraction of the lap elapsed.
  \item \textbf{Pit-stop indicator}: a car entering the pit lane is shown as
        stationary in the pit area for the duration of the stop.
  \item \textbf{Live sidebar}: an always-visible leaderboard lists cars in current
        race order with their accumulated times, updated after every lap.
  \item \textbf{Final podium}: once the race ends and all animations have played
        out, a modal overlay shows the final standings with positions and total times.
\end{itemize}

\begin{figure}[H]
  \centering
  \includegraphics[width=\textwidth]{../../imgs/09_circuit.png}
  \caption{Web dashboard --- Circuit tab: animated Catmull-Rom track with car positions, start lights and live leaderboard sidebar.}
  \label{fig:ui_circuit}
\end{figure}
